% Subsection: Symmetric Difference  (notes stub; content authored later)
\subsection{Symmetric Difference}

\begin{remark*}[Exposition]
Symmetric difference is already a primitive set operation from Volume I.  Its
role here is algebraic: it behaves like addition on subsets.  When subsets of
an ambient set \(X\) are combined by symmetric difference, elements that appear
twice cancel out.
\end{remark*}

\begin{remark*}[Interpretation]
For \(A,B\subseteq X\), the set \(A\triangle B\) contains the points where
membership in \(A\) and membership in \(B\) disagree.  Thus symmetric
difference measures change of membership rather than accumulation of
membership.  This is why it becomes the additive operation in the Boolean-ring
view of set algebras.
\end{remark*}

\begin{remark*}[Algebraic Role]
Inside \(\mathcal{P}(X)\), the empty set acts as the identity for symmetric
difference:
\[
  A\triangle\varnothing=A.
\]
Every subset is its own inverse:
\[
  A\triangle A=\varnothing.
\]
These identities are the set-theoretic source of the characteristic-two
behavior that later appears in Boolean rings.
\end{remark*}

\begin{remark*}[Examples]
If \(A\) and \(B\) are admitted subsets in a family \(\mathcal{F}\), then
\(A\triangle B\) is a subset of the same ambient set \(X\).  Whether it is
again a member of \(\mathcal{F}\) is a closure question.
\end{remark*}

\begin{dependencies}
\begin{itemize}
  \item \hyperref[def:family-of-subsets]{Family of Subsets}
  \item \hyperref[def:sym-diff]{Symmetric Difference}
  \item \hyperref[def:empty-set]{Empty Set}
\end{itemize}
\end{dependencies}
